\documentclass[10pt,a4paper]{amsart}
\usepackage{amsmath,amssymb,amsthm,graphicx}
% Retained minimal preamble for the Stacks Project chapter more-algebra.tex (More on Algebra).
% The chapter file pinned-original-source.tex carries only \input{preamble} on its first line, so no
% preamble text lives inside the chapter itself.  The Stacks Project uses one single shared preamble for
% every chapter; the snapshot stages it identically as sources/stacks-chow/preamble.tex and
% sources/stacks-etale-cohomology/preamble.tex.  The block below is copied mechanically, line for line,
% from sources/stacks-chow/preamble.tex lines 14-18 (the Xy-pic setup) and 154-217 (the native theorem
% declarations together with the mathematical macros).  Nothing was added, removed or redefined; the
% wrapper class and the amsmath/amssymb/amsthm/graphicx base packages are supplied by the assembler.
% Chapter: more-algebra.tex -- More on Algebra.
% For commutative diagrams we use Xy-pic
\usepackage[all]{xy}

% We use 2cell for 2-commutative diagrams.
\xyoption{2cell}

% Theorem environments.
%
\theoremstyle{plain}
\newtheorem{theorem}[subsection]{Theorem}
\newtheorem{proposition}[subsection]{Proposition}
\newtheorem{lemma}[subsection]{Lemma}

\theoremstyle{definition}
\newtheorem{definition}[subsection]{Definition}
\newtheorem{example}[subsection]{Example}
\newtheorem{exercise}[subsection]{Exercise}
\newtheorem{situation}[subsection]{Situation}

\theoremstyle{remark}
\newtheorem{remark}[subsection]{Remark}
\newtheorem{remarks}[subsection]{Remarks}

\numberwithin{equation}{subsection}

% Macros
%
\def\lim{\mathop{\mathrm{lim}}\nolimits}
\def\colim{\mathop{\mathrm{colim}}\nolimits}
\def\Spec{\mathop{\mathrm{Spec}}}
\def\Hom{\mathop{\mathrm{Hom}}\nolimits}
\def\Ext{\mathop{\mathrm{Ext}}\nolimits}
\def\SheafHom{\mathop{\mathcal{H}\!\mathit{om}}\nolimits}
\def\SheafExt{\mathop{\mathcal{E}\!\mathit{xt}}\nolimits}
\def\Sch{\mathit{Sch}}
\def\Mor{\mathop{\mathrm{Mor}}\nolimits}
\def\Ob{\mathop{\mathrm{Ob}}\nolimits}
\def\Sh{\mathop{\mathit{Sh}}\nolimits}
\def\NL{\mathop{N\!L}\nolimits}
\def\CH{\mathop{\mathrm{CH}}\nolimits}
\def\proetale{{pro\text{-}\acute{e}tale}}
\def\etale{{\acute{e}tale}}
\def\QCoh{\mathit{QCoh}}
\def\Ker{\mathop{\mathrm{Ker}}}
\def\Im{\mathop{\mathrm{Im}}}
\def\Coker{\mathop{\mathrm{Coker}}}
\def\Coim{\mathop{\mathrm{Coim}}}

% Boxtimes
%
\DeclareMathSymbol{\boxtimes}{\mathbin}{AMSa}{"02}

%
% Macros for moduli stacks/spaces
%
\def\QCohstack{\mathcal{QC}\!\mathit{oh}}
\def\Cohstack{\mathcal{C}\!\mathit{oh}}
\def\Spacesstack{\mathcal{S}\!\mathit{paces}}
\def\Quotfunctor{\mathrm{Quot}}
\def\Hilbfunctor{\mathrm{Hilb}}
\def\Curvesstack{\mathcal{C}\!\mathit{urves}}
\def\Polarizedstack{\mathcal{P}\!\mathit{olarized}}
\def\Complexesstack{\mathcal{C}\!\mathit{omplexes}}
% \Pic is the operator that assigns to X its picard group, usage \Pic(X)
% \Picardstack_{X/B} denotes the Picard stack of X over B
% \Picardfunctor_{X/B} denotes the Picard functor of X over B
\def\Pic{\mathop{\mathrm{Pic}}\nolimits}
\def\Picardstack{\mathcal{P}\!\mathit{ic}}
\def\Picardfunctor{\mathrm{Pic}}
\def\Deformationcategory{\mathcal{D}\!\mathit{ef}}

% --- appended 2026-10-12 (Chief): the project own definitions of the reference/slogan/history
% display environments, copied mechanically from the shared Stacks preamble lines 8-12
% (\usepackage{verbatim} plus the three \newenvironment lines, which the chapter preamble keeps
% verbatim from the project preamble). Nothing added, removed or redefined: these environments
% carry no mathematical content and the project itself typesets them as comments. Without these
% three lines a counted statement whose own text uses them fails with Environment slogan undefined.
\usepackage{verbatim}
\newenvironment{reference}{\comment}{\endcomment}
\newenvironment{slogan}{\comment}{\endcomment}
\newenvironment{history}{\comment}{\endcomment}
\usepackage{cleveref}
\crefname{originaldefinition}{original definition}{original definitions}
\Crefname{originaldefinition}{Original definition}{Original definitions}
\makeatletter
\newlabel{derived-definition-saturated}{{05RB}{source}{}{sourcepointer-1}{}}
\newlabel{derived-definition-saturated@cref}{{[originaldefinition][0][]05RB}{[1][0][]source}{}{}{}}
\makeatother
\crefname{originaldefinition}{original definition}{original definitions}
\Crefname{originaldefinition}{Original definition}{Original definitions}
\makeatletter
\newlabel{derived-definition-generators}{{09SJ}{source}{}{sourcepointer-2}{}}
\newlabel{derived-definition-generators@cref}{{[originaldefinition][0][]09SJ}{[1][0][]source}{}{}{}}
\makeatother
\crefname{originallemma}{original lemma}{original lemmas}
\Crefname{originallemma}{Original lemma}{Original lemmas}
\makeatletter
\newlabel{lemma-two-out-of-three-perfect}{{066R}{source}{}{sourcepointer-3}{}}
\newlabel{lemma-two-out-of-three-perfect@cref}{{[originallemma][0][]066R}{[1][0][]source}{}{}{}}
\makeatother
\crefname{originallemma}{original lemma}{original lemmas}
\Crefname{originallemma}{Original lemma}{Original lemmas}
\makeatletter
\newlabel{lemma-summands-perfect}{{066S}{source}{}{sourcepointer-4}{}}
\newlabel{lemma-summands-perfect@cref}{{[originallemma][0][]066S}{[1][0][]source}{}{}{}}
\makeatother
\title{The ring itself generates the category of perfect objects as the smallest saturated triangulated subcategory}
\author{The Stacks Project Authors}
\date{stacks-project:more-algebra.tex:local-snapshot-20261012; GNU FDL 1.2 or later}
\begin{document}
\maketitle
\setcounter{subsection}{1}
\noindent\textit{Editorial note.} All delivered mathematical text is the author's own, copied from the pinned Stacks Project chapter file more-algebra.tex; no statement or proof step was rewritten, reordered or removed. The counted statement is the labeled lemma environment at source lines 20287-20293 (canonical tag 0ATI) and the counted proof is the proof environment at source lines 20295-20323, which immediately follows it and is its own proof.\par\medskip
\noindent\textit{Editorial note.} The proof begins with a sentence addressed to the reader, asking to look at two definitions to see what the statement means; that sentence is part of the author's own proof inside the counted span and its references are carried as navigation stubs.\par\medskip
\noindent\textit{Editorial note.} The four references are two to definitions of another chapter, to which the proof explicitly sends the reader, and two to earlier lemmas of this chapter establishing the structural properties; all four are navigation stubs with their roles recorded in the reasons.\par\medskip
\noindent\textit{Editorial note.} The preamble span runs to line 90, the current end of the retained preamble, which includes the appended environment definitions; the counted text does not use them, and the span is taken whole so that the file version is unambiguous.\par\medskip
\noindent\textit{Editorial note.} The source URL is the source-level Stacks Project URL, as required when more than one card is produced from one source; the per-claim tag is carried by claim\_key and by this note. The counted statement and proof cite no bibliography entry, so this card needs no bibliography.\par\medskip
\noindent\textit{Allowed standard prerequisites (editorial).}\begin{itemize}
\item The perfect objects form a strictly full, saturated triangulated subcategory of the derived category. invoked at proof lines 20299-20301. Shows that the subcategory in question contains the smallest such subcategory generated by the ring.
\item The subcategory generated by an object is the union of its finite stages, the first stage containing every finite projective module placed in any degree. used at proof lines 20303-20309. Provides the induction target and the base case of the induction.
\item A complex concentrated in degrees from a to b sits in a distinguished triangle with its top term shifted and its truncation below the top degree. used at proof lines 20315-20322. Realizes the induction step, showing that the complex belongs to the next stage of the generated subcategory.
\end{itemize}\medskip
\ifdefined\hypertarget\hypertarget{sourcepointer-1}{}\fi
\noindent\textit{Source reference.} Original definition 05RB: sources/stacks-chow/tags.txt, source lines 7405--7405. The reference fixes the meaning of a saturated subcategory, invoked at proof lines 20296-20297 so that the statement can be read; it is terminology and the mathematical content is inside the counted span.\par\medskip
\ifdefined\hypertarget\hypertarget{sourcepointer-2}{}\fi
\noindent\textit{Source reference.} Original definition 09SJ: sources/stacks-chow/tags.txt, source lines 12357--12357. The reference fixes the meaning of generation and of the stages of the generated subcategory, invoked at proof lines 20296-20304; the induction that uses the stages is inside the counted span.\par\medskip
\ifdefined\hypertarget\hypertarget{sourcepointer-3}{}\fi
\noindent\textit{Source reference.} Original lemma 066R: sources/stacks-more-algebra/pinned-original-source.tex, source lines 16848--16848. The reference is one of the two cited at proof lines 20299-20300 for the statement that perfect objects form a saturated triangulated subcategory; the induction built on that is inside the counted span.\par\medskip
\ifdefined\hypertarget\hypertarget{sourcepointer-4}{}\fi
\noindent\textit{Source reference.} Original lemma 066S: sources/stacks-more-algebra/pinned-original-source.tex, source lines 18963--18963. The reference is the second of the two cited at proof lines 20299-20300, supplying closedness under direct summands; it is a standard pointer and everything drawn from it is inside the counted span.\par\medskip
\begin{lemma}
\label{lemma-perfect-ring-classical-generator}
Let $R$ be a ring. The full subcategory $D_{perf}(R) \subset D(R)$ of perfect
objects is the smallest strictly full, saturated, triangulated subcategory
containing $R = R[0]$. In other words $D_{perf}(R) = \langle R \rangle$.
In particular, $R$ is a classical generator for $D_{perf}(R)$.
\end{lemma}

\begin{proof}
To see what the statement means, please look at
Derived Categories, Definitions \ref{derived-definition-saturated} and
\ref{derived-definition-generators}.
It was shown in Lemmas \ref{lemma-two-out-of-three-perfect} and
\ref{lemma-summands-perfect} that $D_{perf}(R) \subset D(R)$
is a strictly full, saturated, triangulated subcategory of $D(R)$.
Of course $R \in D_{perf}(R)$.

\medskip\noindent
Recall that $\langle R \rangle = \bigcup \langle R \rangle_n$.
To finish the proof we will show that if $M \in D_{perf}(R)$
is represented by
$$
\ldots \to 0 \to M^a \to M^{a + 1} \to \ldots \to M^b \to 0 \to \ldots
$$
with $M^i$ finite projective, then $M \in \langle R \rangle_{b - a + 1}$.
The proof is by induction on $b - a$.
By definition $\langle R \rangle_1$ contains any finite projective
$R$-module placed in any degree; this deals with the base case
$b - a = 0$ of the induction. In general, we consider the distinguished
triangle
$$
M_b[-b] \to M^\bullet \to \sigma_{\leq b - 1}M^\bullet \to M_b[-b + 1]
$$
By induction the truncated complex $\sigma_{\leq b - 1}M^\bullet$ is
in $\langle R \rangle_{b - a}$ and $M_b[-b]$ is in $\langle R \rangle_1$.
Hence $M^\bullet \in \langle R \rangle_{b - a + 1}$ by definition.
\end{proof}

\end{document}
